\section{Methodology}
\label{sec:methodology}

This section describes the six processing stages of the framework in the
order data flows through them: TLS preprocessing (\ref{sec:preprocess}),
geometric object decomposition with structure-remnant filtering
(\ref{sec:decomposition}), explicit modeling and provisional classification
(\ref{sec:explicit}), multi-view VLM verification with automatic escalation
(\ref{sec:verification}), object identity matching (\ref{sec:matching}), and
confidence-gated graph upsert (\ref{sec:upsert}), followed by the structurally
gated query interface (\ref{sec:query}) and the object-grounded place layer
(\ref{sec:place}). Two further mechanism families---the corrective
decomposition layers with the structure-condition ensemble, and the
owner-evidence revision path---are defined in Sec.~\ref{sec:corrective},
with the failures that motivated them reported in
Secs.~\ref{sec:exp-seg} and~\ref{sec:exp-owner}; for brevity we refer
to the geometric stages through explicit modeling as \emph{Stage~A} and
to VLM verification as \emph{Stage~B}. Stages \ref{sec:preprocess}--\ref{sec:explicit} form the
\emph{deterministic geometric backbone}: every random draw is seeded, so a
re-run on the same host reproduces every retained cluster byte-identically
(Sec.~\ref{sec:exp-seg}). Stages \ref{sec:verification} onward are \emph{probabilistic
semantic inference}, controlled by voting, escalation, and confidence gating
rather than by determinism claims. The pipeline is fully automated
\emph{after registered TLS input}: the only manual steps are operating the
scanner and the vendor registration of multi-setup scans.

\subsection{TLS Preprocessing}
\label{sec:preprocess}

The input is a registered E57 point cloud from a survey-grade terrestrial
scanner (Leica BLK360). When the export contains multiple setups, the loader
applies each setup's registered pose and merges them into a single cloud; a
70M-point six-setup scan and a pre-merged single-scan export are treated
identically. The cloud is voxel-downsampled at 3\,cm---dense enough to retain
object detail while bounding the cost of all later stages. The cloud is kept
in the scanner's registered frame; no global normalization is applied.
Instead, the floor height is estimated where a stage needs it, by a robust
statistic suited to that stage's input (the densest 5\,cm height bin of the
lower half-cloud for the remnant filter; the lower quartile of per-cluster
bottom heights for the map-based filter of
Sec.~\ref{sec:corrective}).

Structural surfaces are then removed by iterative RANSAC plane
extraction~\cite{fischler1981ransac}. Each round draws 300 seeded random
point triples, keeps the candidate with the most inliers at a 5\,cm distance
threshold, and refits it once by a least-squares (SVD) fit to those inliers.
A plane is removed only if it is structural by orientation: near-horizontal
($|n_z|\ge0.85$; floors, ceilings, split ceiling levels) or near-vertical
($|n_z|\le0.2$; walls). The loop stops at the first of four conditions:
(i)~ten planes have been removed, (ii)~fewer than 2{,}000 points remain,
(iii)~the best remaining plane supports less than 4\% of the points still
present---no large plane is left---or (iv)~the best remaining plane is
slanted ($0.2<|n_z|<0.85$), i.e., object-like, so removing it would begin
carving objects. The fixed seed makes the extraction deterministic; the
pre-removal cloud is retained as a \emph{ceiling reference}
for the remnant filter of Sec.~\ref{sec:decomposition} and for incremental
re-runs. Fig.~\ref{fig:pipeline}a,b show a scan before and after
structural-plane removal.

\subsection{Geometric Object Decomposition and Structure-Remnant Filtering}
\label{sec:decomposition}

The structure-free cloud is partitioned into object candidates by
fixed-parameter DBSCAN~\cite{ester1996dbscan}
(\mbox{$\varepsilon=0.30$\,m}, 100-point minimum),
followed by a \emph{conservative giant-split}: any cluster whose horizontal
footprint exceeds 2.5\,m is re-clustered at a tighter
\mbox{$\varepsilon=0.20$\,m}, and the split is accepted only if it yields at
least two sub-clusters of 100+ points---separating merged neighbors without
shattering genuinely large equipment. We position this stage deliberately:
it is not a novel segmentation algorithm but a \emph{reproducible
decomposition designed for verification and incremental diffing}. Its value
is that identical inputs yield identical clusters with stable ordering, which
later stages exploit for result carry-over and change detection
(Sec.~\ref{sec:matching}). Fig.~\ref{fig:pipeline}c shows the resulting
instances, and Fig.~\ref{fig:pipeline}d two of the TOSM records they
eventually populate. The parameter sensitivity of this choice is
quantified in Sec.~\ref{sec:exp-seg}. Four automatic corrective layers and
one human-in-the-loop channel extend this conservative backbone; they are
defined, with their trigger conditions and parameters, in
Sec.~\ref{sec:corrective}.

\paragraph{Structure-remnant filter.}
Real buildings violate the implicit assumption that structure is planar and
objects are not. Our showroom scene has a three-level ceiling
($z=2.48/2.18/1.95$\,m): RANSAC removes each horizontal ceiling plane, but the
\emph{vertical soffit bands} at the level transitions survive removal and
surface as phantom objects (Fig.~\ref{fig:pipeline}e)---thin, wall-colored
slabs that a downstream verifier confidently misreads (e.g., as wall-mounted
keyboards or handrails, Sec.~\ref{sec:exp-verify}). We therefore filter any
cluster that passes all of the following tests, with the thresholds fixed
once and used unchanged on every scene. (i)~\emph{Soffit-scale geometry}: the
cluster's vertical extent is at most 0.6\,m and its bottom is at least
1.5\,m above the floor. (ii)~\emph{Thin vertical planar band}: the best-fit
plane (SVD) has $|n_z|\le0.25$ and an RMS off-plane thickness of at most
3\,cm. (iii)~\emph{Ceiling above}: the cluster's footprint, padded by
0.3\,m, contains at least 40 pre-removal points in the band from 2\,cm to
45\,cm above the cluster top, and (iv)~\emph{dense horizontal layer}: at
least 60\% of those points concentrate in a single $\sim$7\,cm slab around
the modal 3.5\,cm height bin---i.e., the band hangs from a real ceiling.
(v)~\emph{Topmost layer}: the region above that slab (0.1--1.5\,m higher)
holds at most 20\% as many points as the slab itself; window sills and wall
remnants fail this test because wall continues above them. The filter reduced the
showroom scene from 45 to 38 clusters and the robot hall from 88 to 81 while
leaving every retained cluster byte-identical, so previously computed
verification results remain valid. Horizontal and sloped ceiling patches are
out of the filter's scope and are discussed as a residual limitation in
Sec.~\ref{sec:discussion}.

\begin{figure}[!tp]
\centering
\includegraphics[width=0.92\linewidth]{fig2_pipeline.png}
\caption{Geometric backbone of TLS-SMF, illustrated on the test-room scene.
(a)~Registered multi-station TLS scan after preprocessing.
(b)~The same cloud after seeded-RANSAC removal of structural planes (floor,
walls, ceiling).
(c)~Deterministic decomposition of the residual cloud into object instances;
colors denote cluster identity, and identical inputs reproduce identical
clusters with stable ordering.
(d)~Two of the resulting TOSM knowledge-graph records with their symbolic,
explicit, and implicit attributes; \texttt{machine\_004} carries an
owner-verified type correction (\emph{machine}\,$\rightarrow$\,\emph{Mobile
Robot}).
(e)~Structure-remnant filter, shown as a front elevation of the showroom
scene: the vertical soffit bands at the three-level ceiling transitions
(red) survive horizontal-plane removal and would otherwise surface as
phantom objects; the filter removes them while leaving every retained
cluster byte-identical.}
\label{fig:pipeline}
\end{figure}

\subsection{Explicit Modeling and Provisional Classification}
\label{sec:explicit}

Each cluster is summarized by an explicit model: axis extents and centroid,
yaw from the principal component of the horizontal point distribution,
dominant color, and point count. These attributes are deterministic functions
of the cluster; the backbone's determinism claim covers geometry and these
explicit attributes only---the provisional label below is stochastic and
excluded from it. A provisional semantic label is then attached by Uni3D
open-vocabulary classification over a 30-class industrial vocabulary. The
provisional label is explicitly \emph{not trusted}: our own vocabulary
ablation shows that narrowing the vocabulary to the deployment
domain \emph{reduces} accuracy (Sec.~\ref{sec:exp-ablation}). The
provisional label serves only as VLM context and as a fallback name; every
label a robot may act on must pass the verification stage below.

\subsection{Multi-View VLM Verification with Automatic Escalation}
\label{sec:verification}

Verification renders each object from four azimuths ($0^{\circ}$,
$90^{\circ}$, $180^{\circ}$, $270^{\circ}$) and submits each view
independently to a vision--language model under a forced tool-use
schema: the model must return a type, a confidence, and a stated reason, with
the object's measured dimensions and provisional label provided as context.
Free-text answers are impossible by construction, and prompts, decoding
parameters, and model version are recorded with every result. The verifier
model itself is a swappable component, not part of the design: the concrete
model is an experimental control (Sec.~\ref{sec:setup}), and Sec.~\ref{sec:exp-ablation} repeats the
identical protocol across five commercial VLMs to characterize how the
choice affects decisiveness, accuracy, and cost.

Per-view votes are combined by confidence-weighted late fusion. Let view $v$
return label $\ell_v$ (lowercased; no synonym mapping is applied at fusion
time) with confidence $c_v\in[0,1]$. Each candidate label accumulates the
\emph{sum} of the confidences of the views that voted for it,
$s(\ell)=\sum_{v:\,\ell_v=\ell}c_v$; the fused type is
$\hat{\ell}=\arg\max_{\ell}s(\ell)$, and the \emph{vote share} is its
normalized confidence mass, $\sigma=s(\hat{\ell})/\sum_{\ell}s(\ell)$.
A confidence-weighted sum, rather than a view count, lets one hesitant
dissenter matter less than a confident one. Three outcomes follow
(Fig.~\ref{fig:escalation}). If the votes are
unanimous, the object is \texttt{verified}; unanimity is empirically a strong
signal (94.4--100\% precision on the two scenes' confirmed ground truth,
Sec.~\ref{sec:exp-verify}),
although Sec.~\ref{sec:exp-verify} also documents its failure modes. If the vote splits, the
object \emph{escalates automatically}: eight additional images---four diagonal
azimuths and four zoomed close-ups---are classified, and the fusion is
recomputed over all twelve votes pooled (the base four included). If the
post-escalation share reaches $\sigma\ge0.6$
the object becomes \texttt{verified\_escalated}; otherwise it remains
\texttt{unverified}. The 0.6 threshold was fixed a~priori, before any
evaluation, as the smallest round value that demands a clear supermajority
of the confidence mass: at 0.5 a twelve-view vote split between two labels
of equal confidence could verify on a tie-break, whereas 0.6 requires the
winning label to dominate all alternatives combined by a margin. The
threshold gates \emph{status only}---the fused label is recorded either
way---so its effect is a precision--recall trade on the verified tier;
Sec.~\ref{sec:exp-verify} re-scores the stored votes across a threshold
sweep to quantify that trade. Crucially, unresolved objects are not hidden, discarded,
or resolved by a human: they stay in the graph with their failed candidate
label attached, and the query interface (Sec.~\ref{sec:query}) is responsible
for making their status impossible to ignore. Implicit robot-facing
attributes (movability, openability, key-object status) are inferred once at
fusion time from the winning view set.

Escalation bounds cost: only split votes pay for the wider view set
(trigger rates 51--58\%, Sec.~\ref{sec:exp-verify}), and per-call cost is $\sim$\$0.009, so a
scene verifies for a few dollars (Sec.~\ref{sec:discussion}).

\begin{figure}[t]
\centering
\includegraphics[width=0.95\linewidth]{fig3_multiview_escalation.png}
\caption{Multi-view verification and automatic escalation, with per-view
votes, fused share, and the three terminal statuses.}
\label{fig:escalation}
\end{figure}

\subsection{Object Identity Matching}
\label{sec:matching}

The update problem is an \emph{identity} problem before it is a geometry
problem: a robot that plans against node IDs needs yesterday's chair to be
today's chair. Our policy is \textbf{mint-once}: a node ID is minted when an
object first enters the graph---derived from a fingerprint of its birth-time
anchor (map, position, extent)---and never changes afterwards. The
fingerprint is a candidate-retrieval key, not the identity itself; after
minting, matching alone carries identity forward.

Matching between a new scan's records and existing nodes runs in two passes:

\begin{enumerate}
\item \textbf{Spatial pass.} A record matches a node if their 3-D centroids
agree within 0.5\,m and each of the three extents (length, width, height)
agrees within a per-axis tolerance of $\max(0.3\cdot
\max(d_{\mathrm{old}},d_{\mathrm{new}}),\,0.15\,\mathrm{m})$---a 30\%
relative band with a 15\,cm absolute floor so that small objects are not
over-constrained. No type constraint applies in this pass; a changed label
on a geometrically stable object is recorded as an update, not a new
object. Matches update pose and extents in place (\texttt{unchanged} /
\texttt{updated}).
\item \textbf{Move pass.} Remaining records match remaining nodes of the
\emph{same type} with dimensions consistent under the same per-axis
tolerance, within a 6\,m horizontal search radius,
producing \texttt{moved} nodes with preserved IDs. Generic types
(\texttt{clutter}, \texttt{unknown}) are excluded from this pass: a generic
label carries too little identity for displacement claims, and our real
two-epoch data showed that admitting them fabricates 2--5\,m
``moves'' between unrelated piles (Sec.~\ref{sec:exp-update}).
\end{enumerate}

\emph{Conflict resolution.} When several nodes satisfy the predicate for one
record, the record claims the node with the smallest horizontal centroid
distance; when several records compete for one node, the first claim wins
and the node is removed from the candidate pool, so later records fall
through to the move pass and then to insertion. Assignment is therefore
greedy in record order rather than a global optimum---a deliberate
simplicity choice: records arrive in the backbone's deterministic cluster
order, so the assignment itself is reproducible, and at the object densities
of our scenes ($\le$0.6 objects/m$^2$) multi-candidate conflicts are rare
(Sec.~\ref{sec:exp-update} reports the observed conflict counts).
The thresholds encode physical priors rather than tuned values: 0.5\,m is
below the smallest inter-object spacing we expect a scan-to-scan pose drift
to produce for a stationary object (and is far above the $\le$5\,cm
registration residual, Sec.~\ref{sec:setup}); the 30\%/0.15\,m extent band
absorbs partial-view truncation between scans without admitting
differently-sized objects; and 6\,m bounds a within-room relocation. They
were fixed on the development epochs (T1/T2) and used unchanged afterwards.

Records that match nothing are minted as new nodes (\texttt{inserted});
nodes that no record claims are marked \texttt{absent} but retained with
their full history, so an object that returns later resumes its original
identity. Sec.~\ref{sec:exp-update} evaluates matching accuracy, identity-switch rate, and
false-new/false-absent rates, including the occlusion level at which
matching breaks (40\% crop).

\subsection{Confidence-Gated Graph Upsert}
\label{sec:upsert}

Each scan epoch applies its matched record set to the knowledge graph as a
node-level upsert---no rebuild. Every node carries its verification status,
provenance (per-view votes, fused share, escalation flag, model version), a
per-revision history of every change, and \texttt{firstSeen}/\texttt{lastSeen}
timestamps. Edges (spatial relations---\texttt{isNextTo}, \texttt{isOn},
\texttt{isAboveOf}---recomputed from explicit models) are scoped to the map
and revision, and organized by the place layer: \texttt{isNextTo} is kept
only between objects of the same place (proximity across a place boundary
is expressed at the place level as \texttt{isAdjacentTo}), vertical
relations must pass a rotated-footprint overlap test (a separating-axis
check that discards the long-range stacked relations oversized merged
clusters otherwise produce), and each place designates a key object
(the VLM's implicit \texttt{isKeyObject} among verified members, falling
back to the most confident specific-type verified member; objects refuted
by owner ground-truth feedback are excluded from designation while
remaining in the graph as annotated errors). On the T3
epoch this scoping removes 24 of 49 raw proximity edges---long
cross-place spokes centered on the largest cluster---leaving 25
intra-place edges over 11 place adjacencies; Fig.~\ref{fig:relmap} shows
the resulting map. The result is exported both as a canonical
JSON graph and as mapId-scoped Neo4j Cypher, so multiple sites and epochs
coexist in one database. Because unchanged objects reproduce byte-identical
clusters (Sec.~\ref{sec:decomposition}), their verification results are
carried over by content hash and an identical-input repeat scan costs
\emph{zero} VLM calls;
a changed scan of the same acquisition quality pays only for changed
clusters (3.1\% of full cost in our controlled combined-edit measurement,
Sec.~\ref{sec:exp-update}). The carry-over is conditional on reproducible
acquisition: when scan quality itself changes between epochs, as in our
real T2$\rightarrow$T3 transition, no cluster hashes match and the epoch
pays full verification cost (Sec.~\ref{sec:exp-update}).
Fig.~\ref{fig:upsert} summarizes
the mint-once identity flow and the node outcomes of a re-scan.

\begin{figure}[t]
\centering
\includegraphics[width=0.85\linewidth]{fig12_relations_map.png}
\caption{Place-scoped TOSM relation map (T3) over the place layer of
Fig.~\ref{fig:placemap} (pastel): dotted object$\rightarrow$place
\texttt{isInsideOf} containment, intra-place \texttt{isNextTo} edges,
vertical \texttt{isOn}/\texttt{isAboveOf} relations, dashed place
\texttt{isAdjacentTo} edges between region centroids, and per-place key
objects (named stars). The map reflects the graph's current revision,
including owner feedback (Sec.~\ref{sec:exp-owner}).}
\label{fig:relmap}
\end{figure}

\begin{figure}[t]
\centering
\includegraphics[width=0.9\linewidth]{fig5_identity_upsert.png}
\caption{Identity matching and confidence-gated upsert: mint-once IDs,
two-pass matching, and the five node outcomes across a re-scan.}
\label{fig:upsert}
\end{figure}

\paragraph{Formal summary.}
Consolidating the above: an object record is the tuple
$r=(a, p, d, \ell, c, \tau, \pi)$ with birth anchor $a$ (map ID plus
position and extent quantized at 0.25\,m, hashed once at first insertion),
pose $p\in SE(2)\times\mathbb{R}$, extents $d=(l,w,h)$, fused label $\ell$,
confidence $c$, verification tier $\tau$, and provenance $\pi$ (per-view
votes, share, escalation flag, model version). The tier set is
\[\begin{aligned}\mathcal{T}=\{&\texttt{verified},\ \texttt{verified\_escalated},\ \texttt{verified\_majority},\\
&\texttt{verified\_recovery},\ \texttt{verified\_owner},\ \texttt{unverified},\ \texttt{refuted}\},\end{aligned}\]
ordered by evidence
source: unanimity, escalated supermajority, agreement-gated recovery, and
owner declaration (confidence 1.0); the query gate
(Sec.~\ref{sec:query}) exposes only the verified tiers as typed objects.
The match function is
$m(n,r)=\mathbb{1}[\lVert p_n-p_r\rVert_{3\mathrm{D}}\le0.5]\cdot
\prod_{k\in\{l,w,h\}}\mathbb{1}[\,|d_{n,k}-d_{r,k}|\le
\max(0.3\max(d_{n,k},d_{r,k}),0.15)]$, with the move-pass variant
substituting a 6\,m horizontal bound and an exact-type, non-generic
constraint. The upsert decision for each record--node pair follows the
deterministic cascade \emph{match} $\rightarrow$ \emph{move} $\rightarrow$
\emph{insert}, with unclaimed nodes marked \emph{absent}; every transition
appends to the node's history, and no node is ever deleted. All numeric
constants above and in
Secs.~\ref{sec:preprocess}--\ref{sec:upsert} were fixed on the development
epochs (T1/T2) and applied unchanged to T3 and the cross-floor scene.

\subsection{Corrective Decomposition Layers and Owner Feedback}
\label{sec:corrective}

The backbone of Sec.~\ref{sec:decomposition} is deliberately conservative,
and its characteristic failures---diffuse residual webs, wall-flush
compounds amputated by structure removal, structure noise that survives
plane removal, and real objects absorbed into \texttt{clutter}---are
corrected by five layers defined here. The first four are automatic and
run inside the frozen pipeline on every epoch; the fifth is the
human-in-the-loop channel. Each was \emph{developed} in response to
failures observed on the development epochs
(Sec.~\ref{sec:exp-owner} preserves that discovery narrative and its
evaluation); their definitions, triggers, and parameters were frozen
before the cafeteria evaluation and require no per-scene parameter
adjustment at execution; held-out transfer is evaluated only where
explicitly reported in Sec.~\ref{sec:exp-endtoend}.

\begin{enumerate}
\item \textbf{Diffuseness-gated re-split.} Trigger: a cluster whose 2-D
occupancy fill (10\,cm grid) is below 0.30 while its footprint exceeds
5\,m$^2$---real objects in our scenes fill 0.5--0.85. The cluster is
re-decomposed by occupancy-core erosion: grid cells with $\ge$8 points and
$\ge$3 occupied 8-neighbors seed connected components; components with
$\ge$250 points survive. Compounds are further split at the support plane
(the $z$-histogram peak 0.5--0.95\,m above floor), with grounded cells
spanning $>$1.3\,m vertically treated as legs/poles, and under-slab
fragments re-attached at $\ge$60\% footprint overlap. An evidence gate
($\ge$250 points, $\ge$5\,m$^2$ or 12\,m$^2$ for desk banks,
$\ge$0.15\,m height) keeps only physically plausible sub-objects.
\item \textbf{Wall-compound gate.} Trigger: $\ge$25\% of a cluster's
points within 0.12\,m of the exploration-grid wall line over a span of
$\ge$2\,m, with $\ge$250 points protruding off-wall, and the cluster not
diffuse. The wall band is stripped, the support plane is detected as
above, and the panel band is cut at density valleys along the wall,
yielding unit-level decomposition of wall-mounted equipment. A final
\emph{panel square-up} pass then re-assigns every compound point that
falls inside a panel unit's bounding box, padded by 6\,cm horizontally
and 5\,cm vertically, back to that unit---drawing from the extracted
support-post pieces and from points left unassigned by the cuts---so that
post extraction and layer clustering do not leave bite marks or holes in
the panel rectangles; it is instance-segmentation finishing with no
semantic input.
\item \textbf{Structure-condition ensemble.} Trigger: every epoch.
The geometric pass is repeated under three structure-removal conditions
(none / ceiling-only / full); clusters found under an alternative
condition but absent from the production set become new candidates, and
only those are sent to verification. Cross-condition label agreement is
stored per object as a \texttt{labelStability} attribute
(stable/unstable) that prioritizes owner review; cross-condition
\emph{voting} is deliberately not used to fuse labels, because verifier
biases proved correlated across conditions
(Sec.~\ref{sec:exp-seg}).
\item \textbf{Hybrid grid-map structure filter.} Trigger: every epoch,
after clustering. Each cluster is registered against the exploration
occupancy grid (an independent, drift-free wall prior) and flagged by the
first matching rule of a fixed cascade: outside-room ($>$50\% of points
beyond the filled room polygon), ceiling fixture (bottom $>$2.0\,m above
floor), wall remnant (wall-hug fraction---the fraction of cluster points
within 0.15\,m of the nearest grid wall line---$>$0.6 with PCA minor
extent $<$0.22\,m and a horizontal minor axis), or full-height wall band
($z$-span $>$2.5\,m, per-metre-segment median thickness $<$0.3\,m,
along-wall length $>$2\,m---the length condition spares wall-mounted
panels). Flagged clusters are removed upstream of verification.
\item \textbf{Owner feedback and gated recovery.} The console exposes
four operations, each applied as a normal KG revision with history:
\emph{declare} (sets \texttt{verified\_owner} at confidence 1.0),
\emph{refute} (sets \texttt{refuted}, presence \texttt{absent}, and
exclusion from key-object designation), \emph{correct} (label or
attribute in place, the VLM verdict retained in provenance), and
\emph{re-verify}. Recovery re-verification re-asks the verifier with
\texttt{clutter} treated as \emph{unidentified}, in two independent
passes (open vocabulary; site vocabulary plus Uni3D candidates); a
recovered label is adopted---as \texttt{verified\_recovery}---only when
the two passes agree under canonical synonym matching, or a pass agrees
with the cluster's own base label, and the label is not \texttt{clutter};
a focused tie-break pass re-asks with exactly the two disputed labels and
is accepted only when its answer is one of them at confidence $\ge$0.6.
\end{enumerate}

\subsection{Query Interface with Structural Gating}
\label{sec:query}

The graph is queried through an LLM that receives a serialized view of the
relevant nodes. The naive design---instructing the model not to assert
unverified labels---fails in practice: in our benchmark the instructed model
still asserted trap labels in 3 of 4 cases (Sec.~\ref{sec:exp-ablation}). We therefore
\emph{gate by construction, not exhortation}: in the query-time view, every
unverified node's type is demoted to \texttt{unknown}, and the failed
candidate label is moved to a separate \texttt{failedCandidateType} field
whose semantics (``a classifier guess that did not survive verification'') are
stated in the schema, not in behavioral instructions. Fig.~\ref{fig:cards}
contrasts the two record states as the query interface sees them. With the
gate in the data rather than the prompt, the model cannot assert an
unverified label without visibly contradicting its input: the four
unverified-label traps are blocked in the gating-style ablation (4/4,
Sec.~\ref{sec:exp-ablation}), while in the end-to-end benchmark one count-phrased trap still
leaks through the quarantine field itself (3/4, Sec.~\ref{sec:exp-query}). Traps whose
label was wrongly \emph{verified} lie beyond any status-based gate---a
bound we quantify in Sec.~\ref{sec:exp-query}.

In interaction, the gate makes confidence part of the answer rather than
metadata. Asked ``is there a keyboard in the room?'', the system answers
from the gated view: a verified keyboard is asserted with its location,
while an unverified one yields ``an object at $(x,y)$ was proposed as a
keyboard by the classifier but did not pass verification''---an answer a
robot can act on conservatively (route a person, lower approach speed)
instead of a confident guess.

\begin{figure}[t]
\centering
\includegraphics[width=0.95\linewidth]{fig4_record_cards.png}
\caption{TOSM semantic-object records (symbolic / explicit / implicit
layers). (a)~A verified record asserts its type. (b)~An unverified record
is structurally gated at query time: the type is demoted to
\texttt{unknown} and the failed classifier guess is quarantined in
\texttt{failedCandidateType}.}
\label{fig:cards}
\end{figure}

\subsection{Object-Grounded Place Layer}
\label{sec:place}

Robot tasking refers to places (``the docking area'') as often as to objects.
Following the TOSM place schema of DK-SMF~\cite{joo2025dksmf}---boundary
polygon, containment, and purpose---we add a place layer grounded in the
verified object layer. Place regions are obtained automatically from the
robot-facing occupancy grid map (0.05\,m---generated from the scan here,
and equally the exploration SLAM map on deployment), following DK-SMF's
recipe: the room interior delimited by the mapped outer boundary is
partitioned by SLIC superpixel segmentation~\cite{achanta2012slic} with the
number of superpixels as the governing parameter (7 requested here),
fragments below an area threshold (1.5\,m$^2$) are discarded, and each
region's centroid is the moment-based center of its cells (cf. grid-map
room segmentation~\cite{bormann2016room,luperto2022robust}). Each
object's \texttt{isInsideOf} resolves to the region containing its anchor
cell. Each place is then summarized by its
\emph{ring code}: the sequence of member objects sorted by bearing around the
place centroid---restricted to verified-tier objects, so the confidence gate
propagates into place semantics---and a single LLM call names the place from
its ring code. Because the segmentation depends only on the map, region
identity is stable across scan epochs, and the same regions re-derive names
that track the room's actual use (the machine bay: exhibition-era
\texttt{machine\_access\_area} $\rightarrow$ storage-era
\texttt{clutter\_zone} $\rightarrow$ test-era
\texttt{machine\_operator\_station}; Sec.~\ref{sec:exp-update}), demonstrating that place
semantics stay current without manual re-annotation.
Fig.~\ref{fig:placemap} shows the resulting T3 place layer. The robot
layer of TOSM is out of scope here.

\begin{figure}[t]
\centering
\includegraphics[width=0.85\linewidth]{fig11_place_map.png}
\caption{Object-grounded place layer at T3. SLIC superpixel
regions~\cite{achanta2012slic} tile the room interior delimited by the
nav-stack grid map's outer boundary and carry LLM-derived names computed
from verified-object ring codes; unverified objects (blue markers) do not
contribute, so the confidence gate propagates into place semantics. Shown
at the graph's owner-corrected revision (Sec.~\ref{sec:exp-owner}).}
\label{fig:placemap}
\end{figure}
